\section*{Philadelphia-Hähnchen}
\ihead{}\chead{Personen:3}\ohead{}
\ifoot{Vorbereitungszeit:20 min}\cfoot{}\ofoot{Kochzeit:40 min}
{\Large Zutaten}
\begin{multicols}{2}
\begin{itemize}
\item \SI{500}{g} Hähnchenbrustfilet(s) oder Putenbrustfilet
\item \SI{150}{g} Kräuterfrischkäse, z.,B. Philadelphia
\item \SI{200}{ml} Sahne
\item \SI{100}{ml} Milch
\item \SI{1}{TL} Brühepulver nach Wahl
\item Schinkenspeck, geräuchert, in Scheiben
\item Senf, mittelscharf
\item Gartenkräuter
\item evtl. Salz und Pfeffer
\end{itemize}
% \columnbreak
\end{multicols}
\noindent {\Large Zubereitung}\\
\\
Jede Filethälfte einmal teilen und großzügig mit Senf bestreichen.
Die Fleischstücke in Schinkenspeck wickeln und in eine feuerfeste Form legen.
Im auf \SI{160}{\degreeCelsius} Umluft vorgeheizten Ofen ca. \SI{25}{min} braten.
Bei Ober-/Unterhitze den Ofen auf \SI{180}{\degreeCelsius} vorheizen.
Inzwischen Sahne, Milch, Brühepulver, Frischkäse und ggf. gehackte frische Gartenkräuter, z.,B. Petersilie, Kresse, Kerbel oder Schnittlauch, zu einer Soße verrühren.
Die Soße über die Filets geben und die Form für weitere \SI{15}{min} in den Ofen geben.
Salz und Pfeffer sind nicht unbedingt nötig, da Brühe, Käse, Senf und Schinkenspeck bereits für ausreichend Würze sorgen.
Dazu passt Reis, besonders gut Curryreis.
